\documentclass{article}
\title{Fractions and Roots}
\author{ }
\date{}

\begin{document}
\maketitle

\section*{1. Fractions}
A fraction has a numerator (top part) and a denominator (bottom part).

\textbf{LaTeX source code:}
\begin{verbatim}
\[
\frac{a}{b}
\]
\end{verbatim}
\textbf{Expected compiled output:}
\[
\frac{a}{b}
\]
\textbf{Explanation:} \verb|\frac| creates a fraction. The first brace group is the numerator and the second is the denominator.

\section*{2. Square roots}
\textbf{LaTeX source code:}
\begin{verbatim}
\[
\sqrt{x}
\]
\end{verbatim}
\textbf{Expected compiled output:}
\[
\sqrt{x}
\]
\textbf{Explanation:} \verb|\sqrt| creates a square-root symbol.

\section*{3. Other roots}
\textbf{LaTeX source code:}
\begin{verbatim}
\[
\sqrt[3]{x}
\]
\end{verbatim}
\textbf{Expected compiled output:}
\[
\sqrt[3]{x}
\]
\textbf{Explanation:} The optional value in square brackets gives the root index. Here, 3 means cube root.

\section*{4. Practice}
Write one-half, the square root of \(a+b\), and the fourth root of \(x\).
\end{document}
